\section{Historie}
Prolínání her a umělé inteligence (AI) je dlouhodobě zkoumanou oblastí. Velká pozornost ve výzkumu AI týkající se her se zaměřuje na vývoj agentů schopných zapojit se do hry, ať už s integrací strojového učení, nebo bez ní. Tato cesta zkoumání historicky představovala primární a často jediný přístup k využití umělé inteligence v herním kontextu.

Ještě před formalizací umělé inteligence jako samostatného oboru se průkopníci počítačové vědy snažili prozkoumat možnosti výpočetní techniky tím, že navrhovali programy schopné hrát stolní hry. Mezi těmito prvními snahami je práce Alana Turinga, který je obecně považován za zásadní postavu počítačové vědy a který upravil algoritmus Minimax pro hraní šachů.

Počátky zapojení umělé inteligence do digitálních her lze vysledovat až k průkopnické práci A.S. Douglase z roku 1952, kdy v rámci své doktorské práce na univerzitě v Cambridge vyvinul software schopný naučit se dokonale hrát Tic-Tac-Toe. Příspěvky Arthura Samuela následně znamenaly významný milník v podobě zavedení takzvaného \uv{reinforcement} učení, předchůdce současných technik strojového učení, demonstrovaného na programu, který autonomně zlepšoval svojí taktiku v dámě prostřednictvím opakovaného samohraní\cite{Samuel1959}.

Raný výzkum v oblasti umělé inteligence zaměřené na hraní her se převážně soustředil na tradiční deskové hry, jako je dáma nebo šachy. Tyto hry byly obzvláště zajímavé díky své schopnosti generovat složité komplexy z jednoduchých souborů pravidel a díky svému historickému postavení po staletí testovala hranice lidského poznání. V roce 1994 došlo k významnému průlomu, když program Chinook Checkers porazil úřadujícího mistra světa v dámě Mariona Tinsleyho. Následně v roce 2007 byla dáma zcela vyřešena\cite{SBBKMLLS2007}, což znamenalo významný milník ve výzkumu AI.

\newpage

Šachy, často označované jako \uv{octomilka umělé inteligence}, posloužily jako klíčové měřítko pro testování nových metod umělé inteligence. Vývoj softwaru, který je schopen překonat lidský výkon v šachu, však posunul zaměření výzkumu AI na jiné oblasti. Deep Blue\cite{WikiDB} od společnosti IBM, využívající algoritmus Minimax rozšířený o úpravy specifické pro šachy a vysoce optimalizovanou funkci vyhodnocování šachovnic, v roce 1997 slavně porazil mistra světa Garryho Kasparova.

Dalším významným pokrokem, který předcházel úspěchům Deep Blue a Chinook, je demonstrován softwarem TD-Gammon\cite{WikiTD}, který v roce 1992 vyvinul Gerald Tesauro. TD-Gammon se vyznačuje využitím umělé neuronové sítě vycvičené pomocí \uv{temporal difference learning}, která se sama proti sobě opakovaně zapojuje do hry. Pozoruhodné je, že TD-Gammon dosáhl dovedností, které jsou srovnatelné s uznávanými lidskými hráči vrhcábu, což dokazuje účinnost jeho tréninkové metodiky.

Po průlomových objevech v oblasti umělé inteligence v tradičních deskových hrách dosáhla vrcholu v oblasti umělé inteligence v deskových hrách hra Go v roce 2016. Hra Go se stala novou hračkou pro herní umělou inteligenci, která se vyznačuje bezkonkurenčním faktorem větvení blížícím se k číslu 250 a rozsáhlým prohledávacím prostorem, který řádově převyšuje prostor šachů. Ze začátku bylo dosažení lidské úrovně v Go považováno za vzdálenou budoucnost, ale nakonec v říjnu 2015 odehrál AlphaGo, který patří společnosti Google DeepMind a je vybaven revolučním frameworkem \uv{Deep reinforcement learning}, svůj první zápas proti úřadujícímu trojnásobnému mistru Evropy Fan Hui, který i vyhrál a skóroval 5:0. V roce 2016 Lee Sedol, profesionální hráč Go s devíti tituly, podlehl porážce v zápase na pět her proti AlphaGo. Díky těmto drtivým vítězstvím se hra Go stala posledních stolní hrou, kde počítače dosáhly nadlidských schopností v klasických deskových hrách, protože další stolní hry zatím lidské hráče v širokém měřítku nezaujaly.\cite{AlphaGo}

Tradiční deskové hry, které se vyznačují strukturovaným tahovým hraním a transparentním sdílením informací mezi hráči, však představují pouze podmnožinu herních prostředí. Uznává se, že inteligence zahrnuje širší dimenze než ty, které jsou testovány v těchto tradičních hrách.

V posledních patnácti letech se vytvořila rostoucí komunita, která se zaměřuje především na zlepšování schopností umělé inteligence v digitálních hrách s cílem dosáhnout výkonnosti srovnatelné s lidskými hráči, replikovat specifické lidské strategie nebo vykazovat jedinečné vlastnosti.

K významnému pokroku v této oblasti došlo v roce 2014, kdy algoritmy vyvinuté společností Google DeepMind prokázaly schopnost vyniknout v různých hrách z klasické konzole Atari 2600 a překonaly výkon na úrovni člověka pouze za použití surových pixelových vstupů. Mezi těmito hrami představovala výzvu pro konvenční přístupy hra Ms. Pac-Man (Namco, 1982). Tým Microsoft Maluuba však tuto výzvu vyřešil použitím techniky posilovacího učení s hybridní architekturou odměn, čímž nakonec dosáhl řešení této hry.

\newpage

V poslední době však došlo k prudkému nárůstu výzkumu zaměřeného na vývoj věrohodných agentů ve hrách, což otevřelo cestu novým směrům v herní umělé inteligenci. 

Jednou z cest zkoumání je vytváření agentů schopných projít herním testem, který se jmenuje Turing test\cite{OGDD2021}. Tento test, odvozený od klasického Turingova testu, vyžaduje, aby porota přesně rozeznala, zda pozorované herní chování pochází od lidského hráče nebo od herního bota řízeného umělou inteligencí. V roce 2012 se dva boty řízené umělou inteligencí ve hře Unreal Tournament 2004 od Epic Games úspěšně prokázaly schopnost projít herním Turingovým testem, což svědčí o jejich schopnosti replikovat herní chování podobné lidskému.

\newpage

\section{\uv{Chování pro tento případ} neboli Ad-Hoc Behavior}
V oblasti vývoje her je již dlouho známé tzv. \uv{Ad-Hoc Behavior}, jako jsou konečné stavové stroje, stromy chování a umělá inteligence založená na užitku. Tyto metody tradičně umožňují značnou kontrolu nad nehráčskými postavami (NPC). O jejich trvalé převaze svědčí i to, že se s jejich používáním v komunitě vývojářů her stále spojuje termín \uv{herní umělá inteligence}.

\subsection{Konečný automat (Finite State Machines)}
Až do poloviny 20. století byl dominantní metodou pro ovládání a rozhodování nehráčských postav (NPC) ve hrách konečný automat(Finite State Machines, KA)\cite{Gill207} a jeho varianty, například hierarchické KA. Tyto techniky měly v oblasti herní umělé inteligence převahu a prokázaly svou účinnost při řízení chování NPC v herním prostředí.

Konečné automaty spadají do oblasti expertních znalostních systémů a běžně se vizualizují jako grafy. Tyto grafy slouží jako abstraktní reprezentace vzájemně propojených prvků, jako jsou objekty, symboly, události, akce nebo vlastnosti relevantní pro konkrétní kontext návrhu. V rámci této reprezentace uzly odpovídají stavům, z nichž každý je charakterizován matematickou abstrakcí, zatímco hrany označují přechody, které znamenají podmíněné vztahy mezi stavy. KA fungují tak, že v daném okamžiku může být aktivní pouze jeden stav, přičemž k přechodům dochází na základě splnění předem definovaných podmínek. V podstatě jsou charakterizovány třemi základními prvky:
\begin{enumerate}
	\item Stavy:
	\begin{itemize}
		\item Obsahují informace týkající se úlohy, jako je například aktuální stav \uv{zkoumání}. 
	\end{itemize} 
	\item Přechody:
	\begin{itemize}
		\item Označují změny stavu vyvolané podmínkami, které musí být splněny, jako je přechod do stavu \uv{upozorněn} po zaslechnutí střelby.
	\end{itemize} 
	\item Akce:
	\begin{itemize}
		\item Určují specifické chování, jež má být provedeno, například náhodný pohyb a hledání protivníka ve stavu \uv{průzkum}.
	\end{itemize} 
\end{enumerate}
Konečné automaty nabízejí pozoruhodné výhody díky své jednoduchosti při návrhu, implementaci, vizualizaci a ladění. Při aplikaci na rozsáhlejší systémy však mohou představovat výzvu kvůli své potenciální složitosti, což vede k výpočetním omezením v některých aspektech. 

\newpage

Kromě toho mají významnou nevýhodou, spolu s dalšími ad-hoc behavior metodami, a to je jejich přirozená nedostatečná flexibilita a dynamika, pokud nejsou speciálně navrženy tak, aby tyto vlastnosti zahrnovaly. Jakmile jsou dokončeny, otestovány a odladěny, jejich přizpůsobivost a evoluční kapacita jsou omezeny, což vede k projevům relativně předvídatelného chování v herním kontextu.

\begin{figure}[ht]
\includegraphics[scale=0.8]{1-1.png}
\centering
\caption{Příklad KA pro jednoduchou hru Snake. Obsahuje tři stavy a čtyři přechody. Had se postupně přibližuje k jídlu, ale pokud má před sebou zeď a nebo svoje tělo, tak se těmto prvkům pokusí vyhnout.}
\label{obrKASnake}
\end{figure}

\newpage

\subsection{Stromy chování (Behavior Trees)}
Strom chování (Behavior Tree, SC)\cite{Colledanchise2018} funguje jako expertní znalostní systém, podobně jako konečný automat (Finite State Machine, KA), který usnadňuje přechody mezi konečnou množinou úloh nebo chování. Významná výhoda SC oproti KA spočívá v jejich modulární struktuře, která umožňuje skládání složitých chování z jednodušších úloh. Na rozdíl od KA, které představují především stavy, jsou SC strukturovány kolem chování, čímž se odlišují od hierarchických KA. Podobně jako KA nabízejí SC snadný návrh, testování a ladění, což přispělo k jejich širokému rozšíření při vývoji her po úspěšných implementacích v titulech, jako je Halo 2\cite{Isla2005} a Bioshock.

Využívají hierarchickou stromovou strukturu, která se skládá z kořenového uzlu, nadřazených uzlů a jim odpovídajících podřízených uzlů, z nichž každý představuje specifické chování (viz obrázek \ref{obrSCLogika}). Procházení SC začíná od kořene a aktivuje provádění dvojic rodič-dítě, jak je ve stromu naznačeno. Během předem stanovených časových kroků může podřízený uzel vrátit svému rodiči jednu z následujících hodnot: \uv{běží}, pokud chování zůstává aktivní, \uv{úspěch} při úspěšném dokončení nebo \uv{selhání}, pokud chování selže. SC se skládají ze tří základních typů uzlů: sekvence, selektor a dekorátor, z nichž každý plní odlišné funkce, které jsou podrobněji popsány níže:
\begin{enumerate}
	\item Sekvence (na obr.\ref{obrSCLogika} znázorněna šedou barvou):
	\begin{itemize}
		\item Úspěch v podřízeném chování vyvolá pokračování v sekvenci, což nakonec vede k úspěchu nadřazeného uzlu, pokud všechna podřízená chování uspějí; naopak, sekvence selže, pokud některé podřízené chování selže. 
	\end{itemize} 
	\item Selektor (na obr.\ref{obrSCLogika} znázorněn červenou barvou):
	\begin{itemize}
		\item Uzly selektorů zahrnují dva základní typy: selektory pravděpodobnosti a selektory priority. U pravděpodobnostních selektorů jsou podřízená chování vybírána na základě pravděpodobností předem stanovených tvůrcem SC. Naopak prioritní selektory řadí podřízená chování do seznamu a zkoušejí je postupně. Bez ohledu na typ selektoru vede úspěch v podřízeném chování k úspěchu selektoru. Při neúspěchu v podřízeném chování je vybráno další podřízené chování v pořadí (u prioritních selektorů) nebo selže samotný selektor (u pravděpodobnostních selektorů).
	\end{itemize} 
	\item Dekorátor (na obr.\ref{obrSCLogika} znázorněn oranžovou barvou):
	\begin{itemize}
		\item Uzly dekorátoru rozšiřují složitost a možnosti jednotlivých podřízených chování. Mezi příklady dekorátorů patří určení počtu iterací, které podřízené chování podstoupí, nebo přidělení časového limitu pro jeho dokončení.
	\end{itemize} 
\end{enumerate}

\newpage

Ve srovnání s KA nabízejí větší flexibilitu při navrhování a jsou relativně jednodušší na testování. Nicméně mají podobná omezení, zejména pokud jde o jejich statickou povahu, která omezuje jejich přizpůsobivost a dynamičnost. Zatímco pravděpodobnostní výběrové uzly vnášejí prvek nepředvídatelnosti, jsou snahy o zvýšení adaptability prostřednictvím úprav ve struktuře stromů.

\begin{figure}[ht]
\includegraphics[scale=0.67]{1-2.png}
\centering
\caption{Příklad stromu chování}
\label{obrSCLogika}
\end{figure}

\subsection{\uv{AI založená na užitečnosti} neboli Utility System AI}
V odvětví herní umělé inteligence je veliká absence modularity chování v různých hrách a představuje významnou překážku pro dosažení vysoce kvalitní umělé inteligence. K překonání omezení KA a SC v tomto ohledu je stále oblíbenějším přístupem umělá inteligence založená na užitečnosti (Utility System AI, USAI)\cite{Graham2019}. Tato metoda nabízí prostředky pro efektivnější návrh řídicích a rozhodovacích systémů ve hrách. V USAI je ke každé herní instanci přiřazena funkce užitku, která kvantifikuje její důležitost v kontextu hry. Tyto funkce mohou například hodnotit faktory, jako je blízkost nepřítele nebo zdravotní stav agenta. Na základě zvážení všech dostupných užitných hodnot a potenciálních akcí určí umělá inteligence nejzásadnější možnost, kterou je třeba v daném okamžiku upřednostnit. Tento přístup vychází z principů teorie užitku z ekonomie a zahrnuje návrh funkcí užitku, podobně jako vytváření funkcí příslušnosti v teorii fuzzy množin.

Užitky zahrnují široké spektrum ukazatelů, od objektivních údajů, jako je zdraví nepřítele, až po subjektivní faktory, jako jsou emoce a vnímané hrozby. Tyto různorodé metriky spojené s potenciálními akcemi nebo rozhodnutími lze sloučit do lineárních nebo nelineárních vzorců, které slouží k vedení agentů při rozhodování na základě souhrnného skóre užitečnosti. Na rozdíl od sekvenčních rozhodovacích procesů KA a SC umožňuje USAI současné vyhodnocení všech dostupných možností. Přiřazením užitku každé možnosti a výběrem té s nejvyšším užitkem mohou agenti určit nejvhodnější postup. 

\newpage

Tyto hodnoty užitku podléhají pravidelnému přehodnocování, například přepočty probíhající v předem stanovených intervalech během hry, aby se zajistilo sladění s dynamickým stavem hry.

USAI má oproti jiným technikám ad-hoc chování výrazné výhody. Její modularita umožňuje herním agentům rozhodovat se na základě dynamického souboru faktorů, což podporuje adaptabilitu rozhodovacích procesů. Rozšiřitelnost navíc umožňuje bezproblémovou integraci nových úvah podle potřeby. 

Opakovaná použitelnost metody navíc usnadňuje přenos užitkových komponent mezi rozhodnutími a napříč různými hrami, což zvyšuje efektivitu vývoje. Tento přístup se hojně využívá v různých herních žánrech, přičemž jeho významné implementace byly zaznamenány ve hrách, jako je Red Dead Redemption pro výběr zbraní a dialogů, a také ve hře F.E.A.R. pro dynamické taktické rozhodování\cite{Dill2010}.

\begin{figure}[ht]
\includegraphics[scale=0.67]{1-3.png}
\centering
\caption{Řízení chování hry Pac-Man založené na USAI. Hrozba (osa x) je
funkce, která leží mezi 0 a 1 a je založena na aktuální pozici duchů. Pac-Man bere v úvahu aktuální úroveň ohrožení, přiřazuje hodnoty užitku (prostřednictvím tří různých křivek) a rozhodne se sledovat chování s nejvyšší hodnotou užitku. V tomto příkladu roste exponenciálně útěk před duchy s ohledem na ohrožení. Naopak útok na duchy exponenciálně klesá s ohledem na ohrožení. Nakonec hledání a sběr peněz roste lineárně až do určité úrovně ohrožení, od tohoto bodu klesá exponenciálně.}
\label{obrUSAI}
\end{figure}

\newpage

\section{Příklady dalších metod}
\subsection{Procházení stromu}
Procházení stromem\cite{WikiTS} je základní úlohou v teorii algoritmizace a zahrnuje systematické zkoumání všech uzlů v rámci datové stromové struktury. Vzhledem k analogii stromu s grafem je procházení stromů úzce spjato s teorií grafů. Metody procházení stromů se běžně demonstrují na binárních stromech, ale lze je snadno rozšířit i na jiné stromové struktury. Dva základní přístupy, prohledávání do hloubky (DFS) a prohledávání do šířky (BFS), se liší pořadím návštěv uzlů a použitím pomocných datových struktur. DFS, často rekurzivní, obvykle využívá zásobník, a to buď explicitně, nebo implicitně prostřednictvím použití zásobníku volání. Naproti tomu BFS typicky využívá frontu. 

Ve scénářích, kde je cílem cílené vyhledávání hodnot spíše než úplné procházení stromu, umožňují specializované konstrukční techniky, jako je vytváření vyhledávacího stromu (typicky binárního vyhledávacího stromu), efektivní použití přizpůsobených variant algoritmů. Další specializované algoritmy jsou použitelné v problémech prohledávání stavového prostoru, které lze reprezentovat jako stromy. Pozoruhodné je, že heuristické algoritmy, včetně stromového vyhledávání Monte Carlo odvozeného z metody Monte Carlo, našly úspěch v aplikacích umělé inteligence v rámci počítačových her.

\subsection{Evoluční algoritmy}
Evoluční algoritmy (EA)\cite{Yu2010} jsou podmnožinou evolučních výpočtů a představují metaheuristické optimalizační algoritmy založené na populaci. Tyto algoritmy čerpají inspiraci z biologické evoluce a využívají mechanismy, jako je reprodukce, mutace, rekombinace a selekce. V EA představují kandidáti na řešení optimalizačních problémů jedince v populaci, jejichž kvalita se hodnotí pomocí fitness funkce. Iterativním použitím těchto operátorů se populace vyvíjí. Výpočetní složitost, zejména spojená s vyhodnocováním fitness funkcí, často představuje v reálných implementacích EA značnou výzvu. K řešení tohoto problému se používají metody, jako je aproximace fitness. Navzdory své zdánlivé jednoduchosti vykazují EA účinnost při řešení složitých problémů, což naznačuje, že složitost algoritmu nemusí vždy přímo korelovat se složitostí problému.

\newpage

\section{AI ve 3D akčních hrách}
\subsection{Obecné úkoly AI}
V 3D akčních hrách by se měla umělá inteligence umět pohybovat virtuálním světem, takže například ovládat základní pohyby po mapě ale také pokročilé strategie jako například vylézt po zdi nebo skrčit se a projít ventilační šachtou a tak dále. Klíčové jsou také poznávací schopnosti a přizpůsobivost, neboli například dokázat předvídat hráčovi interakce s objekty nebo se samotnou umělou inteligencí. AI by mělo také vynikat v boji a používat různé taktiky, jako je obcházení hráče a nebo koordinované útoky. Kromě boje by měli zvládat i úkoly, jako je interakce s prostředím nebo pomáhat hráčovi dosáhnout cíle. Pokud hra obsahuje i příběhovou stránku, tak by měla AI umět mimikovat emoce, aby dokázala vyjadřovat řadu pocitů, aby hráč se vcítil do své postavy a tím pádem do herního světa.

\subsection{F.E.A.R.}
Hra F.E.A.R.\cite{DD1}\cite{DD2}, zkratka pro First Encounter Assault Recon, debutovala v roce 2005 jako střílečka z pohledu první osoby vyvinutá společností Monolith Productions. Navzdory svému stáří je stále chválena jako ukázkový příklad implementace umělé inteligence v bojových videohrách. Klíč k jejímu trvalému úspěchu spočívá v kombinaci inovativních technik umělé inteligence a strategického herního designu.

Jako první do hry implementovali techniku umělé inteligence zvanou automatizované plánování. Tento přístup, inspirovaný systémem STRIPS\cite{FikesWinter1971} vyvinutým v roce 1971 pro programování robotů, umožňuje nehráčským postavám (NPC) zapojit se do sofistikovaných rozhodovacích procesů v herním prostředí. Konkrétně jsou NPC vybaveny schopností plánovat svůj pohyb, vybírat krytí a rozhodovat o bojových akcích, jako je střelba nebo použití granátu. Tato integrace automatizovaného plánování zahrnuje přiřazení konkrétních cílů každé NPC a modelování možných akcí, které mají v herním světě k dispozici. Analýzou těchto možností mohou NPC vytvářet optimalizované plány pro efektivní dosažení svých cílů.

Každá akce umělé inteligence ve F.E.A.R. se řídí předpoklady a účinky. Předpoklady představují podmínky, které jsou ve světě hry nezbytné pro provedení akce, zatímco účinky označují výsledky vyplývající z provedení akce. Například, pokud je postava pozorována, jak běží k objektu a sráží jej na zem, aby jej použila jako úkryt, znamená to sekvenci plánovaných akcí: pohyb, manipulace s objektem a nakonec hledání úkrytu. To znamená, že před zahájením sekvence akcí je třeba zajistit, aby byl objekt vzpřímený, dosažitelný a schopný poskytnout úkryt. Dodržování předpokladů zabraňuje nelogickým akcím, jako je pokus o manipulaci s objektem, který je již převrácený, nebo pohyb k nedosažitelnému objektu, čímž se zvyšuje realističnost NPC ve hře.

\newpage

Kromě plánovacích funkcí využívá systém F.E.A.R. ke koordinaci provádění formulovaných strategií konečný automat (KA). Programátoři v Monolithu si uvědomily, že mohou mapovat každou akci umělé inteligence ve hře na odpovídající animaci, ať už jde o střelbu ze zbraně, přemisťování nebo interakci s prvky prostředí. Díky tomuto se zjednodušuje složitost chování AI tím, že v rámci KA rozděluje akce do tří různých stavů: pohyb na určené pozice, animační sekvence a specializované interakce s určenými \uv{inteligentními objekty}. Bez ohledu na složitost taktického plánu umělé inteligence jej lze vždy zahrnout do těchto tří stavů KA, což zajišťuje zjednodušené provádění a udržuje soudržnost akcí NPC v průběhu hry.

\begin{figure}[ht]
\includegraphics[scale=0.58]{1-4.png}
\centering
\caption{Konečný automat hry F.E.A.R.\cite{Orkin2015}}
\label{obrFEAR}
\end{figure}

Dalším aspektem designu AI, kterou vývojáři implementovali, je kdy využít AI k dosažení požadovaných výsledků a kdy použít důmyslná řešení a úpravy, neboli \uv{padělat} chytrost AI. Například zatímco jednotliví vojáci ve hře F.E.A.R. nemají povědomí o existenci ostatních, hra vytváří iluzi týmové spolupráce pomocí chytré manipulace. Přidělením samostatných cílů různým postavám, jejichž kombinace simuluje spolupráci, hra naaranžuje scénáře, v nichž například jedna postava poskytuje krycí palbu, zatímco druhá se postaví do boku hráče, aby ho mohli jednodušeji zneškodnit. Ačkoli jsou tyto akce prováděny nezávisle na sobě, jejich synchronizace vytváří dojem koordinovaného manévru. Pro posílení této iluze hra obsahuje dialogy mezi postavami, přestože se navzájem neslyší a neodpovídají si, což přidává další vrstvu realismu a umocňuje hráčův dojem, že bojuje proti koordinovanému týmu.

\subsection{ME: Shadow of Mordor a ME: Shadow of War}
Systém Nemesis\cite{GMT2021}, představený ve hře Middle-Earth: 
Shadow of Mordor v roce 2014, který stále patří mezi nejvyspělejší AI ve videohrách. V podstatě funguje jako dynamický generátor padouchů, který vytváří unikátní protivníky s odlišnými vlastnostmi, osobnostmi a rysy. Tito protivníci hráče zaujmou posměšnými průpovídkami a přispívají k příběhu, který se vyvíjí na základě interakcí s hráčem.

V Middle-Earth: Shadow of War je systém Nemesis druhé generace, který rozšiřuje systém z předchozí hry v mnoha ohledech. Nejenže generuje padouchy, ale také spřádá složité příběhy založené na vztazích hráče s těmito orky. Tento systém důmyslně využívá základní hráčské akce, jako je boj, útěk nebo shromažďování informací, jako stavební kameny příběhu a vytváří tak dynamický příběhový zážitek v rámci akční hry v reálném čase.

Klíčem k úspěchu v systému Nemesis je vytvoření zapamatovatelných postav se silnou osobností a jedinečnými vlastnostmi. Díky dialogům, jako jsou například urážky mířené na hráčskou postavu nebo zpívání, ale i díky vizuálním podnětům se každý ork stává osobitým a zapamatovatelným.

\begin{figure}[ht]
\includegraphics[scale=0.50]{1-5.png}
\centering
\caption{Příklad pevnosti ve hře Shadow of War. Každý ork je unikátní jak vzhledově i chováním. Orkové s modrou lebkou jsou na straně hráče, všichni ostatní jsou nepřátelé. Tento obrázek ukazuje i hiearchii orků, kde ork na nejvyšší pozici je \uv{Overlord}, prostřední pozice je \uv{Warchief} a orkové před hradbami jsou kapitáni, ale také jsou ve hře orkové, co nemají žádné postavení a jsou to buďto vojáci nebo otroci \cite{Parkin2017}}
\label{obrME}
\end{figure}

\newpage

Genialita tohoto systému také spočívá v jeho schopnosti prodloužit tyto příběhy za hranice jednotlivých střetnutí a nabídnout alternativní výsledky, jako je ústup, zrada nebo dokonce vzkříšení poražených orků. To zajišťuje, že příběhy přetrvávají a vyvíjejí se na základě akcí a rozhodnutí hráče. Mezi dalšími rozšířeními těchto příběhů jsou vzácné události, jako je zrada \uv{Warchief} nebo proměna. 

Pro podporu těchto vyvíjejících se příběhů hra také obsahuje hiearchii orkské společnosti, která umožňuje povyšování, degradování a posuny moci mezi orky. Navíc se tento systém dynamicky přizpůsobuje na základě výkonu hráče a přináší výzvy nebo příležitosti k udržení příběhového napětí. Systém Nemesis prosperuje na základě zapojení a interpretace hráče, což umožňuje vznik různých příběhů, díky kterým vznikají individuálních zkušenosti a herní prožitky.